\chapter{工程领域应用案例}
\chapterart{ch06}{本体落地的四个工程领域}

\begin{introbox}
\textbf{【导读】}
理论与语言已经齐备：第 2 章给了知识表示的分层骨架，第 3 章给了从需求走到模型的方法，
第 4 章给了 OWL 与 SPARQL 这套语言工具，第 5 章让推理器真正转了起来。
本章把这套装备带进四个真实领域——智能制造调度、自动驾驶场景理解、
建筑信息模型（BIM）合规检查、航空航天故障模式与影响分析（FMEA）。
读案例时请留意它们的共同模式：\textbf{本体定义约束 + 推理器自动检查}——
这正是第 1 章「知识熵减」在不同行业的兑现。
四个案例彼此独立、可按兴趣选读；但 6.1 的模式总论与 6.6 的跨案例分析建议精读，
它们是把案例经验搬回你自己领域的桥梁。
\end{introbox}

\section{应用模式总论：本体落地的四种方式}

前五章一路走来，我们始终在回答「本体是什么、怎么建、怎么推理」。
从本章起问题换了方向：\textbf{本体在真实工程里究竟扮演什么角色？}
如果把读过的论文与工业报告摊开来看，会发现表象千差万别——
有的项目在车间里检测排产冲突，有的在车载芯片上理解路口场景，
有的对着建筑模型逐条核查防火规范，有的在替新机型推演故障后果。
但剥掉行业外衣之后，本体真正承担的工程职能可以归纳为四种\textbf{应用模式}（application pattern）：
约束检查、语义集成、决策支持、合规审计。
本节先把四种模式各自的通用结构与适用判据讲清楚，
后面四个案例你就能一眼认出「这是哪种模式、为什么选它」。

\subsection{模式一：约束检查——公理与推理器当质检员}

\textbf{约束检查}（constraint checking）是最常见、也最容易见效的模式：
把领域里「什么是不允许的」写成公理与规则，让推理器对源源不断的业务数据做自动质检。

它的通用结构可以用三层文字图来描述：
最上层是\textbf{知识层}——TBox 公理刻画概念边界（如「设备与材料不相交」），
规则库刻画跨实例的违例条件（如「同一设备同一时段不得执行两个任务」）；
中间是\textbf{数据层}——业务系统持续把订单、状态、事件写入 ABox；
最下层是\textbf{服务层}——推理器周期性或事件触发地运行，
凡是被分类进「冲突」「违例」类的个体，连同触发它的公理一起输出成清单。
整个结构像一条质检流水线：知识层是质检标准，数据层是待检品，推理器是那位从不疲倦的质检员。

适用判据有三条，满足越多越值得用这种模式：
\begin{itemize}
\item \textbf{规则可明示}：违例条件能被领域专家用「如果……就不行」的句式说清楚，
  而不是只可意会的手感；
\item \textbf{违例代价高}：漏检一次的损失（停机、返工、事故）远大于建模与维护成本；
\item \textbf{数据已结构化}：待检对象在信息系统里有据可查，
  不需要先做大规模的数据治理。
\end{itemize}

\subsection{模式二：语义集成——给异构系统对齐词汇}

\textbf{语义集成}（semantic integration）解决的是另一类痛苦：
同一个事物在不同系统里叫不同的名字、用不同的结构存储。
设计部门叫「构件」，施工系统叫「工程对象」，运维平台叫「资产」；
字段对不上，含义也对不上，每对接一对系统就要写一套转换脚本，
\(n\) 个系统两两对接要维护接近 \(n^2\) 条管道。

这种模式的通用结构是\textbf{枢纽—辐条}（hub-and-spoke）：
中心是一个共享本体，承载领域的公共词汇与公理；
每个业务系统的局部模式通过一组映射声明挂接到共享本体上
（「系统 A 的 asset 表对应本体的 Equipment 类」）；
查询与交换一律使用共享词汇进行，由映射层负责翻译。
于是 \(n\) 个系统只需维护 \(n\) 组映射，而且映射本身是声明式的、可检查的知识，
不再是埋在脚本里的隐式约定。

适用判据：数据源多于两个且仍在增加；各源词汇差异大、又确实描述同一现实对象；
跨系统查询或跨阶段交付是日常需求而非偶发任务。
若只有两个系统、接口稳定不变，写一段定制转换代码反而更省事。

\subsection{模式三：决策支持——让推荐自带依据}

\textbf{决策支持}（decision support）模式的目标不是替人拍板，
而是把候选方案、约束与偏好显式建模，让推理器给出\textbf{带推理链的建议}。
与黑箱的机器学习推荐不同，本体推理的每一步都援引某条公理或规则，
「为什么推荐它」可以一路追到知识源头。

通用结构是一条四段管道：
\textbf{候选生成}（从 ABox 里枚举可行对象，例如所有空闲设备）
\(\rightarrow\) \textbf{硬约束过滤}（用公理排除不合格者，例如不能加工该材料的设备）
\(\rightarrow\) \textbf{偏好排序}（用规则或属性值对剩余候选打分排序）
\(\rightarrow\) \textbf{解释生成}（把每个候选通过或被排除时触发的公理串成依据链）。
前两段是经典的本体推理，后两段往往与常规程序逻辑配合完成——
本体不必包打天下，它负责的是把「资格判断」这一知识密集环节声明化。

适用判据：决策依据主要是\textbf{知识}（规则、经验、规范）而非海量历史数据的统计规律；
使用者要求解释，「系统说了算但说不出为什么」无法被接受；
错误决策可被追责，需要留下推理依据备查。

\subsection{模式四：合规审计——每条结论可溯源到条款}

\textbf{合规审计}（compliance auditing）可以看作约束检查的加强版，
加强在「溯源」二字：规则不仅要能判违例，还要携带\textbf{出处元数据}——
它形式化自哪部规范的哪一条款、哪个版本、由谁在何时翻译核对。

通用结构：左侧是\textbf{条款库}，每条形式化规则通过注解属性挂着原文条款编号与文本；
中间是\textbf{批量检查引擎}，对整个模型或数据集跑全部规则；
右侧是\textbf{审计报告}，每条违例记录三元组「对象——违反的条款——条款原文」，
监管者或审查师拿着报告可以直接翻回规范原文复核。
与模式一的差别在于产出物：约束检查服务于内部质量，报告给工程师看；
合规审计服务于外部监管，报告要经得起第三方质询，
因此可溯源性从「锦上添花」变成「硬性要求」。

适用判据：领域受强监管（建筑、航空、医疗、金融）；
规范条款数量大、更新频繁，人工核查既慢又容易漏；
审查活动本身需要留痕，「检查过没有、依据是什么」要可证明。

\begin{center}
\begin{tabularx}{\linewidth}{@{}lXX@{}}
\toprule
\textbf{模式} & \textbf{核心机制与典型输出} & \textbf{适用判据} \\
\midrule
约束检查 & 公理与规则定义违例条件，推理器输出冲突清单 & 规则可明示；违例代价高；数据已结构化 \\
语义集成 & 共享本体为枢纽，各系统经映射对齐词汇 & 数据源多且异构；跨系统查询是日常需求 \\
决策支持 & 候选过滤加偏好排序，结论附推理链 & 依据是知识而非统计；必须可解释、可追责 \\
合规审计 & 规则携带条款出处，批量检查生成可溯源报告 & 强监管领域；条款多且常更新；审查需留痕 \\
\bottomrule
\end{tabularx}
\end{center}

\begin{noteblock}
\textbf{模式可以叠加。}
真实项目很少只用一种模式：制造调度以约束检查为主、决策支持为辅；
BIM 项目先做语义集成（IFC 映射）再做合规审计；
FMEA 项目同时是决策支持（影响推演）与合规审计（适航证据）。
选型时先认主导模式——它决定第一期交付什么；
辅助模式放进路线图，等数据与团队成熟再叠加。
\end{noteblock}

带着这张模式地图，我们来预览四个案例。
\textbf{智能制造调度}：把设备能力、任务需求、时间窗口建成本体，
排产冲突由 SWRL 规则自动检出——能力不匹配、状态冲突、时间重叠，
原本靠调度员肉眼核对的工作变成推理器的例行输出。
\textbf{自动驾驶}：车辆要「理解」场景——前方是行人还是锥桶？它和斑马线什么关系？
本体提供交通参与者的语义骨架，感知结果填入成为场景图（scene graph），
每 100 毫秒更新一次；规则在其上判断「行人靠近斑马线 \(\rightarrow\) 准备让行」。
\textbf{BIM 合规检查}：IFC 标准映射为本体后，
设计规范从 PDF 文本变成可执行的 SPARQL 检查——
人工两天的图纸审查缩短到三十分钟，且覆盖全部条款。
\textbf{航空航天 FMEA}：故障模式、影响传播链、风险优先级本体化之后可复用、可推理——
新型号继承既有故障知识库，分析周期从六个月压缩到六周，故障识别量反而提升四成。

\begin{center}
\begin{tabularx}{\linewidth}{@{}lXX@{}}
\toprule
\textbf{领域} & \textbf{主导模式与辅助模式} & \textbf{本体带来的改变} \\
\midrule
制造调度 & 约束检查为主，决策支持为辅 & 冲突检测自动化，调度依据可解释、可审计 \\
自动驾驶 & 约束检查（实时），衔接决策支持 & 场景语义实时化（100ms 级场景图更新） \\
BIM 合规 & 语义集成打底，合规审计落地 & 人工 2 天 \(\rightarrow\) 自动 30 分钟，条款 100\% 覆盖 \\
航空 FMEA & 决策支持与合规审计并重 & 周期 6 个月 \(\rightarrow\) 6 周，故障识别量 +40\% \\
\bottomrule
\end{tabularx}
\end{center}

四个案例的领域差异极大，但工程套路一致：
先用第 3 章的方法划定范围建模，再用第 4 章的语言落地，
最后让第 5 章的推理机制持续工作。
接下来逐案深入，每个案例都按「业务痛点 \(\rightarrow\) 本体设计 \(\rightarrow\) 推理机制 \(\rightarrow\) 价值与教训」展开。

\section{智能制造调度：让排产冲突无处可藏}

\subsection{排产难在哪}

想象一个典型的机加工车间：几十台设备——数控车床、加工中心、磨床、三坐标测量机，
每台的能力参数各不相同；上百张工单在制，每张工单拆成若干道工序，
工序之间有严格的先后约束；材料从铝合金到钛合金，
不是每台设备都能加工每种材料；设备还会进入维护、故障、换型等状态。
排产要回答的问题看似朴素：\textbf{哪道工序、放到哪台设备、在什么时间执行？}

难点不在单个约束，而在约束的\textbf{组合爆炸}与\textbf{持续变化}。
能力约束（这台车床车不了钛合金）、状态约束（那台铣床正在保养）、
时间约束（同一设备不能同时干两件事）、顺序约束（先粗车再精车）彼此纠缠；
而插单、设备故障、来料延迟每天都在发生，昨天排好的计划今天就要重排。
更隐蔽的麻烦是：这些约束大量存在于\textbf{老调度员的头脑}里——
哪台设备「名义上能干、实际上干不好」，哪种材料要避开哪台旧机床，
全是多年踩坑换来的经验。人在知识在，人走知识走，
这正是第 1 章描述的知识熵增场景。

传统软件不是没有尝试：MES 记录了数据但不做约束判断；
APS 排程引擎算力强劲，可是约束以代码形式硬编码在系统里，
工艺一变就要找供应商改程序，周期以月计。
车间需要的是一种\textbf{让约束以知识形态存在}的方案——约束可以被领域工程师直接阅读、
修改、增删，而判断工作交给通用的推理器。这正是本体的用武之地。

\subsection{知识流架构：从需求到推理的五步}

这套调度本体的总体架构，仓库示例文件用一张知识流图概括。

\begin{codebox}
\codeline{\textcolor{cmtgray}{\#\ ==============================}}
\codeline{\textcolor{cmtgray}{\#\ 系统架构（知识流）}}
\codeline{\textcolor{cmtgray}{\#\ ==============================}}
\codeline{\textcolor{cmtgray}{\#}}
\codeline{\textcolor{cmtgray}{\#\ 需求分析\ \&\ CQ\ 能力问题\ \(\rightarrow\)\ 定义\ TBox/RBox/ABox\ 范围}}
\codeline{\textcolor{cmtgray}{\#\ 本体概念建模\ \&\ 类层次设计\ \(\rightarrow\)\ 构建\ TBox}}
\codeline{\textcolor{cmtgray}{\#\ 属性、关系、约束、SWRL\ 规则\ \(\rightarrow\)\ 构建\ RBox}}
\codeline{\textcolor{cmtgray}{\#\ 设备、工艺、工单、任务实例\ \(\rightarrow\)\ 填充\ ABox}}
\codeline{\textcolor{cmtgray}{\#\ 推理、查询、冲突检测\ \(\rightarrow\)\ TBox+RBox\ 推导\ ABox}}
\codeline{\textcolor{cmtgray}{\#}}
\codeline{\textcolor{cmtgray}{\#\ 核心模式：TBox概念\ +\ RBox规则\ \(\rightarrow\)\ 推导ABox隐含知识}}
\codeblank{}
\end{codebox}

\snipsource{ch06-applications/examples/manufacturing-scheduling.txt}

五个箭头串起了一条完整的工程路径，值得逐步拆解。
第一步\textbf{需求分析与能力问题}：用第 3 章的 CQ 方法圈定范围——
本体要回答哪些问题，决定了 TBox、RBox、ABox 各自要装什么，
这一步划下的边界防止后面无限制地膨胀。
第二步\textbf{概念建模与类层次}：把设备、工艺、任务等核心概念整理成 TBox 的类树。
第三步\textbf{属性、约束与规则}：对象属性、属性特征（如传递性）与 SWRL 规则进入 RBox，
这是「领域怎么运转」的知识。
第四步\textbf{实例填充}：来自 MES 与设备台账的具体设备、工单、任务流入 ABox——
注意 ABox 是四步中唯一持续变动的部分。
第五步\textbf{推理与查询}：推理器拿 TBox 和 RBox 当「法律」，对 ABox 里的「案情」做推导。

最后一行的「核心模式」是整章的钥匙：
\textbf{TBox 概念 + RBox 规则 \(\rightarrow\) 推导 ABox 隐含知识}。
显式录入的只有设备参数与任务安排，
而「这个任务有能力冲突」「这两个任务时间重叠」是推理器替你算出来的隐含事实。
对照第 2 章的三层划分可以发现，分层在这里兑现成了\textbf{分工}：
TBox 与 RBox 由本体工程师和工艺专家共同维护、按月演化，
ABox 由业务系统自动写入、按秒变化——稳定的知识与流动的数据各居其位。

\subsection{能力问题设计：范围由问题决定}

第 3 章反复强调：先有能力问题（competency question，CQ），后有本体。
这个案例的五条核心 CQ 如下。

\begin{codebox}
\codeline{\textcolor{cmtgray}{\#\ ==============================}}
\codeline{\textcolor{cmtgray}{\#\ 关键能力问题（Competency\ Questions）}}
\codeline{\textcolor{cmtgray}{\#\ ==============================}}
\codeline{\textcolor{cmtgray}{\#\ CQ1:\ 哪些设备可以加工铝合金？}}
\codeline{\textcolor{cmtgray}{\#\ CQ2:\ 当前有哪些设备空闲可以接受任务？}}
\codeline{\textcolor{cmtgray}{\#\ CQ3:\ 哪些任务存在能力冲突？}}
\codeline{\textcolor{cmtgray}{\#\ CQ4:\ 给定产品，推荐最佳工艺路线？}}
\codeline{\textcolor{cmtgray}{\#\ CQ5:\ 某工单中工序的正确执行顺序？}}
\codeblank{}
\end{codebox}
\color{black}

\snipsource{ch06-applications/examples/manufacturing-scheduling.txt}

逐条审视，每个 CQ 都在向本体「点菜」，规定了它必须具备的构件：
\begin{itemize}
\item \textbf{CQ1（哪些设备可以加工铝合金）}要求存在设备类、材料类，
  以及连接二者的能力关系——这直接催生了对象属性 \texttt{canProcess}；
\item \textbf{CQ2（哪些设备空闲）}要求设备携带状态信息，
  于是有了状态属性 \texttt{hasStatus} 与状态值的枚举；
\item \textbf{CQ3（哪些任务存在能力冲突）}是典型的约束检查需求，
  它要求定义「冲突」这个派生类，并配备产生该类实例的 SWRL 规则；
\item \textbf{CQ4（推荐最佳工艺路线）}把模式从约束检查推向决策支持，
  要求产品、工艺、工序之间的关联可查询、可比较；
\item \textbf{CQ5（工序的正确执行顺序）}要求工序间的先后关系可表达且可传递推导，
  于是 \texttt{precedes} 被声明为传递属性。
\end{itemize}

CQ 在项目里扮演双重角色：建模时是\textbf{范围闸门}——凡是与五条 CQ 无关的概念
（比如设备采购价格、操作工排班）一律不进本体；
验收时是\textbf{测试用例}——本体建成后把每条 CQ 翻译成 SPARQL 查询跑一遍，
答得出、答得对，本体才算合格。这套「CQ 即验收」的做法会在第 9 章的实战中再次出现。

\subsection{TBox 与 RBox 的设计决策}

类层次的全貌在仓库文件 \texttt{manufacturing-scheduling.txt} 中有完整描述：
顶层并列着设备（Equipment）、工艺（Process）、产品（Product）、
材料（Material）、任务（Task）、工单（WorkOrder）六大类；
设备向下按职能分为加工设备、测量设备、辅助设备、工装设备，
加工设备再细分出数控机床、手动机床等；工艺向下挂工序（Operation），
工序再按切削方式分出车削、铣削、钻孔等叶子类。
这棵树上有三个值得说道的设计决策。

\textbf{决策一：设备按职能分类，而不是按品牌或型号。}
品牌型号是实例层面的属性（写进 ABox 的数据属性即可），
职能才决定一台设备「能回答什么 CQ」。
若按型号建类，每采购一台新设备就要改 TBox——
分类轴选错，稳定层与流动层就会互相污染。

\textbf{决策二：工艺与工序分两层。}
工艺（Process）是面向产品的整体方案，工序（Operation）是落到设备上的最小执行单元，
二者以 \texttt{hasOperation} 关联且保持顺序。
这样 CQ4 的「推荐工艺路线」在工艺层回答，CQ5 的「执行顺序」在工序层回答，互不纠缠。

\textbf{决策三：任务与工序分离。}
工序是知识（「精车外圆」这件事该怎么干），任务是执行（今天 14 点在 3 号车床干这件事）。
把静态工艺知识与动态执行安排分开，前者复用率高、后者吞吐量大，
这也是 TBox/ABox 分工原则在业务概念上的投影。

RBox 一侧，核心对象属性同样取自该文件：\texttt{canProcess}（设备到材料，多对多）、
\texttt{hasCapability}（设备到能力）、\texttt{requires}（工序到设备）、
\texttt{uses}（工序到刀具）、\texttt{precedes}（工序到工序，声明为传递属性）、
\texttt{assignedTo}（任务到设备，多对一）。
其中 \texttt{precedes} 的传递性声明最能体现「让推理器干活」的思路：
工艺员只需录入相邻工序的先后（粗车 \texttt{precedes} 精车、精车 \texttt{precedes} 磨削），
「粗车必须在磨削之前」由第 5 章讲过的属性链推理自动补全，
工序再多也不必手工维护全序关系。

\subsection{SWRL 冲突规则精讲}

调度本体的「牙齿」是三条 SWRL 冲突检测规则——第 5 章学过的规则语法在这里上岗。

\begin{codebox}
\codeline{\textcolor{cmtgray}{\#\ ==============================}}
\codeline{\textcolor{cmtgray}{\#\ SWRL调度规则}}
\codeline{\textcolor{cmtgray}{\#\ ==============================}}
\codeblank{}
\codeline{\textcolor{cmtgray}{\#\ 能力匹配：任务需要的材料，设备必须能加工}}
\codeline{Task(?t)\ \(\land\)\ requiresMaterial(?t,\ ?m)\ \(\land\)}
\codeline{assignedTo(?t,\ ?e)\ \(\land\)\ Equipment(?e)\ \(\land\)}
\codeline{\(\lnot\)canProcess(?e,\ ?m)}
\codeline{\hspace*{4\ttcw}\(\rightarrow\)\ CapabilityConflict(?t)}
\codeblank{}
\codeline{\textcolor{cmtgray}{\#\ 状态检查：正在运行的设备不能分配新任务}}
\codeline{Equipment(?e)\ \(\land\)\ hasStatus(?e,\ Running)\ \(\land\)}
\codeline{Task(?t)\ \(\land\)\ assignedTo(?t,\ ?e)}
\codeline{\hspace*{4\ttcw}\(\rightarrow\)\ StatusConflict(?t)}
\codeblank{}
\codeline{\textcolor{cmtgray}{\#\ 时间冲突：同一设备同一时段不能执行两个任务}}
\codeline{Task(?t1)\ \(\land\)\ Task(?t2)\ \(\land\)}
\codeline{assignedTo(?t1,\ ?e)\ \(\land\)\ assignedTo(?t2,\ ?e)\ \(\land\)}
\codeline{hasStartTime(?t1,\ ?s1)\ \(\land\)\ hasEndTime(?t1,\ ?e1)\ \(\land\)}
\codeline{hasStartTime(?t2,\ ?s2)\ \(\land\)\ hasEndTime(?t2,\ ?e2)\ \(\land\)}
\codeline{swrlb:lessThan(?s1,\ ?e2)\ \(\land\)}
\codeline{swrlb:greaterThan(?e1,\ ?s2)}
\codeline{\hspace*{4\ttcw}\(\rightarrow\)\ TimeConflict(?t1,\ ?t2)}
\codeblank{}
\end{codebox}

\snipsource{ch06-applications/examples/manufacturing-scheduling.txt}

逐条解读它们的业务含义与技术细节。

\textbf{能力冲突规则}说的是一句车间常识：「任务要加工的材料，分到的设备必须能干。」
规则体依次绑定任务 \texttt{?t}、它需要的材料 \texttt{?m}、它被指派的设备 \texttt{?e}，
最后一个原子要求设备\textbf{不能}加工该材料，结论把任务归入 \texttt{CapabilityConflict} 类。
注意这里的否定原子在工程上要格外小心：第 5 章讲过 OWL 的开放世界假设——
「知识库里没说能加工」不等于「不能加工」。
实践中的做法是把 \texttt{canProcess} 的完整清单当作局部封闭的数据维护
（设备能力表由台账系统全量同步），在查询层用「找不到能力断言即视为冲突」的封闭式判断实现，
或者显式维护 \texttt{cannotProcess} 断言。
这是把教科书规则落到产线时最典型的一处「翻译损耗」，值得在你自己的项目里提前预案。

\textbf{状态冲突规则}更直白：状态为「运行中」的设备不该再被指派新任务。
它把设备状态（来自数据采集系统的实时信号）与任务指派（来自排产系统的计划数据）
放进同一条规则里比对——两类数据源在传统架构中分属两个系统、靠人对屏幕核对，
而在本体里它们是同一张图上的两条边，推理器顺手就能交叉检验。
语义集成的好处在这条最简单的规则里已经显形。

\textbf{时间冲突规则}最有技术含量：同一设备上的两个任务，
若时间区间相互重叠则报冲突。规则用两个内建比较原子表达重叠——
任务一的开始时间早于任务二的结束时间，且任务一的结束时间晚于任务二的开始时间。
这对「双不等式」正是第 5 章 Allen 区间关系一节的结论：
两个区间存在公共内部，当且仅当各自的开始都早于对方的结束。
看似简单的两行比较，背后是把十三种区间关系里「会打架」的情形一网打尽的完备判定。
工程上还有一个细节：规则对 \texttt{?t1} 与 \texttt{?t2} 是对称的，
同一对冲突会以两种绑定顺序各报一次，落地时要在结果层去重。

三条规则覆盖了排产冲突的三大来源——能力、状态、时间，
而它们的共同结构值得抄录在案：\textbf{规则体描述「坏组合」的特征，
规则头给违例对象贴上冲突类标签}。
冲突类是 TBox 里专设的「告警容器」，下游系统只需订阅这些类的成员变化，
就能驱动看板红灯、消息推送或自动重排。
推理器不修改任何排产决定——它只负责把问题暴露出来，决定权仍在调度员手里。
这种「检测与决策分离」的克制设计，是约束检查模式能被一线接受的关键。

\subsection{应用价值：数字与机制}

按照仓库文件记载的量化结果，这套系统上线后有三项可观测的收益。
其一，\textbf{工艺推荐命中率从零起步，六个月后达到 87\%}——
随着工艺知识持续沉淀进本体，CQ4 那条决策支持线越来越准，
新产品来了先问本体「有没有相似产品的工艺路线可以参考」，老师傅的经验开始可继承。
其二，\textbf{调度冲突检测完全自动化，响应时间小于 100 毫秒}——
插单或改派的瞬间，三条规则同步给出冲突判定，
调度员从「排完祈祷别出错」变成「错了立刻被拦下」。
其三，\textbf{知识复用率上升}：跨产品的工艺推荐减少了重复设计，
同类零件不再每次从零编工艺。

回到 6.1 的模式语言：这是一个以\textbf{约束检查}为主导、
\textbf{决策支持}为辅助的经典组合——冲突检测先上线、先见效、先赢得信任，
工艺推荐随知识积累逐步增值。第 9 章我们会以这个领域为蓝本，
从 CQ 清单开始一行行把系统重建出来，此处先记住它的骨架。

\subsection{经验教训}

复盘这个案例，有三条教训值得带走。

\textbf{教训一：ABox 的实时性比 TBox 的精致度更早成为瓶颈。}
团队前期把大量精力花在类层次的优雅上，
而真正拖慢上线的是设备状态数据的接入——状态不准，
状态冲突规则就会误报漏报，一线很快失去信任。
本体项目的关键路径常常不在本体本身，而在喂给它的数据管道。

\textbf{教训二：规则宁少勿滥。}
每条 SWRL 规则都是一份长期维护负债：工艺变了要改，误报了要查，性能慢了要优化。
三条核心规则覆盖了绝大多数冲突场景，
后续新增规则坚持「先有真实冲突案例、再写规则」的纪律，避免规则库膨胀成无人敢动的黑箱。

\textbf{教训三：CQ 是防过度建模的疫苗。}
项目中途曾有人提议把刀具寿命、夹具库存全部建进本体，
被一个朴素的问题挡了回去——「哪条 CQ 需要它？」
没有问题牵引的概念只会增加维护面而不产生答案。
范围蔓延是本体项目最常见的死法，第 3 章的方法论在这里救了场。

\section{自动驾驶场景理解：感知之上的语义层}

\subsection{感知之后，为什么还需要语义}

现代自动驾驶的感知栈已经强大：神经网络从摄像头与激光雷达数据里
检测出一个个目标——类别、位置、速度、边界框，每秒刷新数十次。
但「检测到」距离「理解了」还差一层。
规划模块真正需要知道的是\textbf{关系与情境}：
那个行人是站在路边等待，还是正在横穿？前方车辆与我同车道吗？
它在减速是因为路口红灯，还是因为更前方有障碍？
这些问题的答案不在任何单个边界框里，而在目标\textbf{之间}、目标与\textbf{道路结构}之间。

纯神经网络方案试图端到端地学出这种情境判断，但它给不出可审查的中间表示——
出了事故无法回答「当时系统认为场景是什么样、为什么这样决策」。
本体走的是\textbf{神经—符号互补}（neuro-symbolic）路线：
感知交给神经网络，它擅长从像素到目标；
语义交给本体，它擅长从目标到情境——用显式的类与关系刻画交通世界的概念结构，
用规则做可解释的风险推断。
感知结果填入本体形成\textbf{场景图}（scene graph）：
节点是交通参与者与道路元素的实例，边是它们之间的空间与语义关系。
场景图是机器的「现场笔录」，每一条边都可以被事后检视。

\subsection{场景图与 100 毫秒更新流程}

场景图不是离线产物，它必须跟上车辆行驶的节奏。仓库文件给出的更新流程如下。

\begin{codebox}
\codeline{\textcolor{cmtgray}{\#\ ==============================}}
\codeline{\textcolor{cmtgray}{\#\ 场景图更新流程（每100ms）}}
\codeline{\textcolor{cmtgray}{\#\ ==============================}}
\codeline{\textcolor{cmtgray}{\#\ 完成对象检测跟踪}}
\codeline{\textcolor{cmtgray}{\#\ \ \ \(\rightarrow\)\ 映射为本体实例（ABox更新）}}
\codeline{\textcolor{cmtgray}{\#\ \ \ \(\rightarrow\)\ 推理空间/时序关系（RBox推理）}}
\codeline{\textcolor{cmtgray}{\#\ \ \ \(\rightarrow\)\ 场景分类（TBox匹配）}}
\codeline{\textcolor{cmtgray}{\#\ \ \ \(\rightarrow\)\ 输出RDF场景图与元数据}}
\codeblank{}
\end{codebox}

\snipsource{ch06-applications/examples/autonomous-driving.txt}

每 100 毫秒——也就是以 10 赫兹的节拍——这条流水线走完四步：
\textbf{第一步}，感知模块完成目标检测与跟踪，输出带跟踪编号的目标列表；
\textbf{第二步}，目标映射为本体实例，写入 ABox——
跟踪编号保证同一个行人在前后帧里是同一个个体，更新它的属性而不是重建它；
\textbf{第三步}，推理空间与时序关系——
根据坐标推出「谁在谁前方」「谁在接近谁」，根据连续帧推出运动趋势；
\textbf{第四步}，场景分类——把当前局面与 TBox 里的场景类型匹配
（路口通行、跟车巡航、行人过街……），输出 RDF 场景图与元数据供下游消费。

把这条流水线与制造案例并排看，会发现\textbf{同构}：
TBox（交通概念体系）在车辆出厂时就固定，
ABox（场景实例）以毫秒级吞吐滚动更新，规则在二者之间持续运转。
不同的只是节拍——车间以分钟计，车载以百毫秒计，
这个量级差异将在 6.3.5 引出真正的工程挑战。

\subsection{核心类与空间关系的设计}

场景本体的类层次在仓库文件 \texttt{autonomous-driving.txt} 中完整列出，骨架是三族概念：
\begin{itemize}
\item \textbf{参与者族}：\texttt{Actor} 为根，下分车辆（\texttt{Vehicle}，
  再细分轿车、卡车、公交车、摩托车）、行人（\texttt{Pedestrian}）、骑行者（\texttt{Cyclist}）。
  分类粒度对齐决策需求——卡车与轿车的让行策略不同，所以值得分；
  轿车的品牌颜色与决策无关，所以不进类树，这与制造案例「按职能不按型号」的决策如出一辙；
\item \textbf{道路族}：\texttt{RoadElement} 下分车道（\texttt{Lane}）、
  路口（\texttt{Intersection}）、人行横道（\texttt{Crosswalk}）——
  它们多来自高精地图而非实时感知，是场景图里相对静态的骨架；
\item \textbf{规制族}：交通标志（\texttt{TrafficSign}）与信号灯（\texttt{TrafficLight}），
  它们把「此处的规则」实体化为场景中的节点。
\end{itemize}

空间关系属性同样取自该文件，分两种风格搭配使用：
\textbf{定性关系}——\texttt{isInFrontOf}、\texttt{isBehind}、\texttt{isLeftOf}、
\texttt{isRightOf}、\texttt{isInSameLane}、\texttt{isApproaching}，
它们把连续的几何世界离散成规则可以直接引用的谓词；
\textbf{定量属性}——\texttt{hasDistance} 保留以米计的精确距离，供需要算术的规则使用。
这种「定性为主、定量兜底」的混合设计是空间建模的常见手法：
定性关系让规则简洁可读，定量属性让阈值判断有据可算。
值得注意的是这些定性关系全部以参照物为中心（多数情况下是自车），
由几何引擎从坐标计算后写入——本体消费几何结论，而不是在本体里做几何运算，
这条「分工线」我们在 BIM 案例还会再遇到。

\subsection{时序规则精讲：从碰撞时间到行人预判}

场景图之上跑着风险推断规则，文件给出两条代表作。

\begin{codebox}
\codeline{\textcolor{cmtgray}{\#\ ==============================}}
\codeline{\textcolor{cmtgray}{\#\ 时序关系推理规则}}
\codeline{\textcolor{cmtgray}{\#\ ==============================}}
\codeblank{}
\codeline{\textcolor{cmtgray}{\#\ 规则：计算碰撞时间（TTC\ -\ Time\ to\ Collision）}}
\codeline{Vehicle(?ego)\ \(\land\)\ Vehicle(?other)\ \(\land\)}
\codeline{isInFrontOf(?other,\ ?ego)\ \(\land\)}
\codeline{hasSpeed(?ego,\ ?v\_ego)\ \(\land\)\ hasSpeed(?other,\ ?v\_other)\ \(\land\)}
\codeline{hasDistance(?ego,\ ?other,\ ?d)\ \(\land\)}
\codeline{swrlb:greaterThan(?v\_ego,\ ?v\_other)\ \(\land\)}
\codeline{swrlb:subtract(?v\_ego,\ ?v\_other,\ ?delta\_v)\ \(\land\)}
\codeline{swrlb:divide(?d,\ ?delta\_v,\ ?ttc)\ \(\land\)}
\codeline{swrlb:lessThan(?ttc,\ 3.0)}
\codeline{\hspace*{4\ttcw}\(\rightarrow\)\ CollisionRisk(?ego,\ ?other)}
\codeblank{}
\codeline{\textcolor{cmtgray}{\#\ 规则：识别行人过街危险场景}}
\codeline{Pedestrian(?p)\ \(\land\)\ isCrossingRoad(?p,\ true)\ \(\land\)}
\codeline{Vehicle(?v)\ \(\land\)\ isApproaching(?v,\ ?p)\ \(\land\)}
\codeline{hasDistance(?v,\ ?p,\ ?d)\ \(\land\)}
\codeline{swrlb:lessThan(?d,\ 30.0)}
\codeline{\hspace*{4\ttcw}\(\rightarrow\)\ PedestrianDangerZone(?p,\ ?v)}
\end{codebox}

\snipsource{ch06-applications/examples/autonomous-driving.txt}

\textbf{第一条是碰撞时间（Time to Collision，TTC）规则}，逐原子拆开看：
前三个原子框定情境——自车 \texttt{?ego}、他车 \texttt{?other}、后者在前者前方；
接着绑定两车速度与车距；
随后三个内建原子做算术与比较——仅当自车更快（\texttt{greaterThan}）才继续，
求出速度差（\texttt{subtract}），再用距离除以速度差（\texttt{divide}）得到 TTC，
即「按当前速度差多少秒后追上前车」；
最后一个原子判断 TTC 小于 3.0 秒，规则头便给这对车辆断言 \texttt{CollisionRisk}。
注意第 5 章强调过的\textbf{绑定顺序纪律}在此体现得淋漓尽致：
\texttt{subtract} 与 \texttt{divide} 的输出变量必须在使用前完成绑定，
而「先判自车更快」不仅是业务语义（更慢就不会追上），
也是技术保护——它排除了速度差为零导致除法无定义的情形。

\textbf{第二条是行人过街危险规则}，它展示了「预判」的符号化做法：
行人 \texttt{?p} 正在横穿马路（\texttt{isCrossingRoad} 为真，
这个判断来自感知层对行人朝向与步态的解读），
某车 \texttt{?v} 正在接近该行人，且车与行人的距离小于 30 米——
三个条件同时成立，断言 \texttt{PedestrianDangerZone}。
规则没有预测行人的精确轨迹（那是概率模型的工作），
它做的是把「值得警惕的局面」定义成可判定的概念：
横穿意图、接近态势、距离阈值，三者都是离散谓词或简单比较，
推理在毫秒级完成，结论带着完整的触发依据。
事后审查时可以精确回答：系统在哪一帧、依据哪三个事实、判定了哪个危险区——
这正是监管与事故分析最需要的可追溯性。

两条规则合起来看，又是熟悉的结构：规则体刻画「危险组合」，
规则头贴上风险类标签——与制造案例的冲突检测完全同构，
只是「冲突」换成了「风险」，节拍从分钟级压到百毫秒级。

\subsection{实时性挑战：推理预算与增量更新}

百毫秒节拍给推理带来了车间里不存在的硬约束：
感知、ABox 更新、规则评估、结果分发全部加起来不能超过 100 毫秒，
留给推理的\textbf{预算}往往只有其中一部分。第 5 章讲过推理复杂度的价目表——
表达力越强代价越高，在车载场景这张价目表变成了生死线。工程上的应对有四招：

\begin{itemize}
\item \textbf{裁剪表达力}：车端本体避开代价高昂的构造
  （大量等价类公理、深层嵌套的存在量词），把 TBox 分类等重活在离线阶段算完，
  车端只做实例级的规则匹配；
\item \textbf{增量推理}：每帧变化的只是少数实例的少数属性，
  只对受影响的规则绑定重新求值，而不是整图重算——
  跟踪编号带来的实例连续性是增量计算的前提；
\item \textbf{规则分级}：碰撞风险这类安全关键规则进高频组、每帧必评，
  场景分类这类辅助规则进低频组、降频运行；
\item \textbf{离线兜底}：一致性检查、不可满足类检测这些昂贵服务放在开发期与
  仿真回放中执行，上车的本体必须是已经验证干净的版本。
\end{itemize}

这套预算思维对所有「本体上实时系统」的场合都适用——
推理不是免费的，表达力与时延的交换必须显式设计，而不是上线后再救火。

\subsection{与规划控制的接口}

语义层不直接踩刹车。它与下游规划控制（planning and control）模块之间是
\textbf{黑板式}的接口：场景图就是那块黑板，
语义层往上写「当前局面是什么、有哪些风险断言」，
规划器从上面读，把 \texttt{CollisionRisk}、\texttt{PedestrianDangerZone}
这类断言当作约束与触发条件——收到行人危险区断言就进入减速让行策略，
收到碰撞风险断言就约束加速度上限。
控制量的连续计算（方向盘转角、制动力度）仍由优化算法完成，
本体负责的是\textbf{把世界翻译成约束}。

这种分工的好处在事故回放时最明显：
黑板上的每一帧场景图都可以存档，
「系统当时看到了什么、推断出什么、为什么这样动作」有完整的符号记录可查。
感知模型可以升级、规划算法可以更换，
而场景图作为中间语言保持稳定——语义层成了整个自动驾驶栈中
最有「合同」性质的一层。

\section{BIM 合规检查：把规范从 PDF 里解放出来}

\subsection{两处断层：模型不连贯，规范不可算}

建筑业的数字化旗舰是\textbf{建筑信息模型}（Building Information Modeling，BIM）：
一栋楼在开工之前先在计算机里被完整地「盖」一遍，
墙、梁、门、管线都是携带属性的三维对象，而不是图纸上的线条。
听起来知识已经高度结构化了，为什么还轮得到本体出场？
因为仓库文件 \texttt{bim-ontology.txt} 开头点出的两处断层依然存在。

\textbf{断层一：跨阶段信息不连贯。}
设计院、施工方、运维方使用不同厂商的软件，
各自的数据模型对「同一堵墙」有不同的编码方式；
设计阶段精心录入的属性，传到施工模型时丢一截，移交运维时再丢一截。
模型在「翻译」中失真，正是 6.1 模式二描述的语义集成问题——
\(n\) 方协作、词汇各异，缺一个公共的语义枢纽。

\textbf{断层二：规范检查不可计算。}
设计要符合的规范——防火、疏散、节能、无障碍——以自然语言条款的形式
躺在规范文本里。审查的方式至今高度人工：
工程师把图纸与条款逐条对照，一个中等规模项目的规范检查要花上\textbf{两天}，
而且条款数以百计，疲劳之下\textbf{遗漏风险}始终存在。
更糟的是设计返工后还要再查一轮，检查成本随设计迭代成倍增长。
这是 6.1 模式四的标准画像：强监管、条款多、人工核查慢且要留痕。

两处断层指向同一条药方：先用本体把模型词汇对齐（语义集成打底），
再把条款写成可执行的规则（合规审计落地）。
仓库文件给出的解决方案正是这条两段式路线：
将 BIM 模型转换为 IfcOWL，建立 BIM 本体定义构件语义，
用 SWRL 规则做规范检查，用 SPARQL 提取合规报告。

\subsection{从 IFC 到本体：语义集成打底}

跨厂商交换的事实标准是 \textbf{IFC}（Industry Foundation Classes），
一个由 buildingSMART 维护的开放数据模式，
定义了上千种建筑实体类型与属性集。
IFC 本身用 EXPRESS 语言书写，社区将其完整翻译成 OWL 本体，
得到的就是 \textbf{IfcOWL}——从此 BIM 模型可以导出为 RDF 图，
墙是个体、楼层是个体、「墙位于楼层」是三元组，
第 4 章的全套查询与推理工具立即可用。

但 IfcOWL 只是「字面翻译」，类型繁多、结构深嵌，直接在其上写规则相当吃力。
工程上的做法是再建一层\textbf{精简的领域本体}，
只保留合规检查关心的概念，向下映射到 IfcOWL。仓库文件给出了它的核心类。

\begin{codebox}
\codeline{\textcolor{cmtgray}{\#\ ==============================}}
\codeline{\textcolor{cmtgray}{\#\ BIM本体核心类（Manchester语法）}}
\codeline{\textcolor{cmtgray}{\#\ ==============================}}
\codeblank{}
\codeline{\textcolor{cmtgray}{\#\ Class:\ BuildingElement}}
\codeline{\textcolor{cmtgray}{\#\ \ \ \ \ SubClassOf:\ owl:Thing}}
\codeline{\textcolor{cmtgray}{\#\ \ \ \ \ Annotations:\ rdfs:label\ "建筑构件"@zh}}
\codeline{\textcolor{cmtgray}{\#}}
\codeline{\textcolor{cmtgray}{\#\ Class:\ Wall}}
\codeline{\textcolor{cmtgray}{\#\ \ \ \ \ SubClassOf:\ BuildingElement}}
\codeline{\textcolor{cmtgray}{\#\ \ \ \ \ Annotations:\ rdfs:label\ "墙体"@zh}}
\codeline{\textcolor{cmtgray}{\#}}
\codeline{\textcolor{cmtgray}{\#\ Class:\ FireCompartment}}
\codeline{\textcolor{cmtgray}{\#\ \ \ \ \ SubClassOf:\ BuildingElement}}
\codeline{\textcolor{cmtgray}{\#\ \ \ \ \ Annotations:\ rdfs:label\ "防火分区"@zh}}
\codeline{\textcolor{cmtgray}{\#}}
\codeline{\textcolor{cmtgray}{\#\ Class:\ ComplianceRule}}
\codeline{\textcolor{cmtgray}{\#\ \ \ \ \ Annotations:\ rdfs:label\ "规范条款"@zh}}
\codeblank{}
\end{codebox}

\snipsource{ch06-applications/examples/bim-ontology.txt}

四个类各有讲究。\texttt{BuildingElement}（建筑构件）是物理对象的根；
\texttt{Wall}（墙体）是典型的可见构件；
\texttt{FireCompartment}（防火分区）值得停一下——它\textbf{不是}模型里画出来的实体，
而是规范视角下的逻辑空间单元，由若干房间与防火墙围合而成；
\texttt{ComplianceRule}（规范条款）则干脆不属于建筑世界，它把「条款」本身实体化，
为的是让违例记录能够指回出处——6.1 说过，可溯源是合规审计的硬性要求。
每个类都带中文 \texttt{rdfs:label} 注解，这是给审查报告的人类读者准备的。

领域类与 IFC 实体的对应关系可以整理成一张映射表：

\begin{center}
\begin{tabularx}{\linewidth}{@{}lXX@{}}
\toprule
\textbf{领域本体类} & \textbf{IFC 侧的对应物} & \textbf{映射要点} \\
\midrule
BuildingElement & IfcBuildingElement 及其子类 & 近乎一对一，直接子类映射 \\
Wall / ExteriorWall & IfcWall，外墙靠属性集与空间关系判定 & 「外墙」在 IFC 中不是独立类型，需推导 \\
FireCompartment & 由 IfcSpace 按防火属性聚合而成 & IFC 无此概念，是规范侧的派生单元 \\
ComplianceRule & 无对应 & 纯规范世界的概念，IFC 不涉及 \\
\bottomrule
\end{tabularx}
\end{center}

这张表暴露了一个深刻的事实：\textbf{模型词汇与规范词汇是两个世界}。
IFC 描述「楼里有什么」，规范谈论「什么样的布置合法」；
「外墙」「防火分区」这些规范概念在模型里没有现成对应，
必须由映射与推理\textbf{派生}出来——比如「外墙」定义为「至少一侧邻接室外空间的墙」，
让推理器把满足条件的 \texttt{IfcWall} 个体自动分类进 \texttt{ExteriorWall}。
本体在这里扮演的正是枢纽角色：左手挂模型词汇，右手挂规范词汇，
把「检查」变成同一张图上的推理。

\subsection{一条条款的完整走查：防火门距离}

方法论说一千道一万，不如把一条条款从纸面送进推理器走完全程。
取一条经简化改写的教学条款：
「\textbf{防火门到最近疏散楼梯间入口的步行距离不应大于 30 米}」
（真实规范按建筑类别与有无自动喷淋分档取值，此处取单一限值只为演示流程）。
形式化分五步。

\textbf{第一步：概念识别。}
条款涉及三个概念——防火门、疏散楼梯间、步行距离，
以及一个判断——距离超限即违例。圈出名词与判断，就得到了建模清单。

\textbf{第二步：词汇对齐。}
查 IFC：「防火门」不是独立实体，模型里只有 \texttt{IfcDoor}，
防火性能藏在属性集的耐火等级字段里。于是在领域本体中定义：
\begin{center}
FireDoor \(\equiv\) Door \(\sqcap\) \(\exists\)hasFireRating.FireRating
\end{center}
「有耐火等级的门就是防火门」——一条等价公理让推理器代替人工去模型里把防火门找齐。
「疏散楼梯间」同理，由楼梯加防火围护的空间组合派生。
这一步最能看出第 3 章概念分析功夫的价值：
条款里的每个名词都要追问「模型里它长什么样」。

\textbf{第三步：可计算化。}
「步行距离」是几何量——沿走道可通行路径的长度，不是两点直线距离。
按 6.3 立下的分工线，几何归几何引擎：
由空间分析程序在模型的通行网络上算出每樘防火门到最近楼梯间的路径长度，
作为数据属性 \texttt{hasDistanceToStair} 写进 ABox。
本体消费几何结论，不在公理里做几何运算。

\textbf{第四步：规则形式化。}写成 SWRL 规则：
\begin{center}
\texttt{FireDoor(?d)} \(\land\) \texttt{hasDistanceToStair(?d, ?dist)} \(\land\)\\
\texttt{swrlb:greaterThan(?dist, 30.0)}
\(\rightarrow\) \texttt{ViolatesEvacuationDistance(?d)}
\end{center}
结构与制造案例的冲突规则完全同构：规则体描述坏组合，规则头贴违例标签。
同时给违例类对应的条款个体挂上注解——
条款编号、原文、版本、翻译人，溯源链就此闭合。

\textbf{第五步：验证。}
拿一个故意做错的测试模型（把某樘防火门挪远）跑规则，
确认违例被检出、报告指向正确条款；再拿合规模型确认无误报。
规则像代码一样需要测试用例——这条纪律在第 8 章评估方法里还会系统展开。

五步走完，一条 PDF 里的条款变成了知识库里\textbf{可执行、可溯源、可回归测试}的资产。
项目里数百条可量化条款，就是沿这条流水线逐一迁移的。

\subsection{规则库与合规查询}

仓库文件给出两条有代表性的正式规则。

\begin{codebox}
\codeline{\textcolor{cmtgray}{\#\ ==============================}}
\codeline{\textcolor{cmtgray}{\#\ 规范自动检查规则（示例）}}
\codeline{\textcolor{cmtgray}{\#\ ==============================}}
\codeblank{}
\codeline{\textcolor{cmtgray}{\#\ 防火规范：防火分区面积不超过2500平方米}}
\codeline{FireCompartment(?fc)\ \(\land\)\ hasArea(?fc,\ ?a)\ \(\land\)}
\codeline{swrlb:greaterThan(?a,\ 2500)}
\codeline{\hspace*{4\ttcw}\(\rightarrow\)\ ViolatesFireCode(?fc)}
\codeblank{}
\codeline{\textcolor{cmtgray}{\#\ 节能规范：外墙传热系数不超过限值}}
\codeline{ExteriorWall(?w)\ \(\land\)\ hasThermalTransmittance(?w,\ ?k)\ \(\land\)}
\codeline{ClimateZone(?zone)\ \(\land\)\ maxTransmittance(?zone,\ ?limit)\ \(\land\)}
\codeline{swrlb:greaterThan(?k,\ ?limit)}
\codeline{\hspace*{4\ttcw}\(\rightarrow\)\ ViolatesEnergyCode(?w)}
\codeblank{}
\end{codebox}
\color{black}

\snipsource{ch06-applications/examples/bim-ontology.txt}

\textbf{防火规范规则}是最简单的「单实体阈值」形态：
防火分区 \texttt{?fc} 的面积 \texttt{?a} 超过 2500 平方米即违例。
别忘了它脚下的地基——\texttt{FireCompartment} 个体是上一小节那类派生推理聚合出来的，
\texttt{hasArea} 是几何引擎算好写入的；规则本身一行，台下功夫十年。

\textbf{节能规范规则}多了一个关键花样：\textbf{参数化条款}。
外墙传热系数的限值不是常数，而是随气候分区变化——
严寒地区一个值、夏热冬冷地区另一个值。
规则没有把限值写死，而是先用 \texttt{ClimateZone(?zone)} 绑定气候分区，
再经 \texttt{maxTransmittance} 原子取出该分区的限值 \texttt{?limit}——
让限值从知识库里\textbf{查出来}。
规范正文后面那张「限值表」被建成了 ABox 断言，
表改了只动数据、不动规则——把「查表」也变成知识，是形式化大型规范的常用手筋。

规则跑完，违例事实已经物化在图里，最后一步是把它们取出来变成审查报告。

\begin{codebox}
\codeline{\textcolor{cmtgray}{\#\ ==============================}}
\codeline{\textcolor{cmtgray}{\#\ SPARQL合规性查询}}
\codeline{\textcolor{cmtgray}{\#\ ==============================}}
\codeblank{}
\codeline{\textcolor{cmtgray}{\#\ PREFIX\ bim:\ <http://example.org/bim\#>}}
\codeline{\textcolor{cmtgray}{\#\ SELECT\ ?element\ ?violation\ ?description}}
\codeline{\textcolor{cmtgray}{\#\ WHERE\ \{}}
\codeline{\textcolor{cmtgray}{\#\ \ \ \ \ ?element\ bim:violates\ ?rule\ .}}
\codeline{\textcolor{cmtgray}{\#\ \ \ \ \ ?rule\ bim:description\ ?description\ .}}
\codeline{\textcolor{cmtgray}{\#\ \ \ \ \ BIND(CONCAT(STR(?element),\ "\ violates\ ",\ STR(?rule))\ AS\ ?violation)}}
\codeline{\textcolor{cmtgray}{\#\ \}}}
\codeline{\textcolor{cmtgray}{\#\ ORDER\ BY\ ?element}}
\end{codebox}
\color{black}

\snipsource{ch06-applications/examples/bim-ontology.txt}

这条 SPARQL 是第 4 章句式的直接复用：
\texttt{WHERE} 里两个三元组模式把「构件——违反的条款——条款描述」连成一行，
\texttt{BIND} 与 \texttt{CONCAT} 拼出人类可读的违例说明，
\texttt{ORDER BY} 让报告按构件排序，便于审查者逐项核对。
注意查询里的 \texttt{bim:violates} 正是 SWRL 规则物化出来的边——
规则负责「判」，查询负责「报」，分工与 6.1 模式四的三段结构严丝合缝：
条款库、检查引擎、审计报告，每条记录都能从报告一路点回条款原文。

\subsection{价值与局限：哪些条款难以形式化}

仓库文件记录的效果有三项：
规范检查从人工 \textbf{2 天}缩短到自动 \textbf{30 分钟}，效率提升 96 倍；
条款覆盖率达到 \textbf{100\%}，机器不会疲劳、不会跳行；
核心类直接映射 IfcOWL，使设计、施工、运维各阶段共享同一套语义，
多规范扩展（再接无障碍规范、绿建标准）只需扩充条款库而不必重建管道。
两天到三十分钟的对比还低估了价值：自动检查可以\textbf{每次设计保存后都跑}，
违例在设计当下被发现，而不是送审前夕——反馈回路从「周」缩到「分钟」，
这才是设计质量提升的真正来源。

跨阶段连贯的价值则要放到建筑的全生命周期里看。
合规检查只是这套语义底座上的第一个应用：
设计阶段沉淀的「这堵墙是防火分区边界」「这樘门是疏散通道上的防火门」，
以 RDF 断言的形式原样流进施工与运维——
施工方据此安排防火封堵的工序验收，
运维方据此生成防火门巡检清单、在改造申请触及防火分区边界时自动告警。
传统流程里这些语义在阶段交接时随格式转换湮灭，运维方往往要对着竣工图重新摸排；
本体让「楼的知识」与「楼」同寿命，这是语义集成模式在时间维度上的延伸——
对齐的不只是并存的系统，还有先后接力的阶段。

但必须诚实地划出边界：\textbf{不是所有条款都能走进知识库}。按形式化难度可分四档：
\begin{itemize}
\item \textbf{量化阈值类}（面积、距离、宽度、系数）：最易形式化，本案例的主力，占比最大；
\item \textbf{需要几何派生的条款}（步行距离、采光角度、视线遮挡）：
  规则不难，难在前置的几何计算，需要空间分析引擎配合；
\item \textbf{语义模糊类}（「宜采用」「重要建筑」「必要时」）：
  条款本身留有解释空间，强行形式化等于替规范编写者做决定，
  工程上只能交回人工判断，或与编制单位确认解释口径后再形式化；
\item \textbf{引用外部知识类}（「按当地主管部门要求」「参照相关标准执行」）：
  指向知识库之外的活动文件，形式化的前提是先把被引用对象本身知识化。
\end{itemize}

成熟项目的做法是\textbf{条款分级混合审查}：
可形式化条款交给机器全量检查，模糊条款保留人工、但由系统自动生成待查清单，
保证「哪些查了机器、哪些查了人、哪些没查」在报告里一目了然。
合规审计模式的可溯源要求，不仅作用于查出的违例，也作用于检查行为本身。

\section{航空航天 FMEA：让故障知识可推理、可复用}

\subsection{表格里的知识困境}

\textbf{故障模式与影响分析}（Failure Mode and Effect Analysis，FMEA）
是安全关键行业的看家功课：对系统里的每个组件，
系统性地追问「它会怎么坏？坏了会怎样？怎么发现？怎么兜底？」
在航空航天领域，FMEA 不是可选的质量活动，而是适航审定的必交材料——
每一个故障模式的影响链都要论证到飞行层面，证明要么后果可接受、要么有缓解措施。

传统 FMEA 的载体是\textbf{表格}：一行一个故障模式，
列依次是组件、故障模式、原因、局部影响、系统影响、飞行影响、探测方法、缓解措施。
表格的问题不在格式而在\textbf{知识形态}：每个单元格都是自然语言短句，
机器读不懂，「液压系统压力丧失」在 A 表里是某行的「系统影响」，
在 B 表里又是另一行的「起因」，同一事物没有同一身份。
由此派生出三大痛点：
\textbf{其一，影响传播靠人脑推演}——泵堵塞导致什么、再导致什么，
全凭分析师在脑中走链路，链条长了就漏；
\textbf{其二，设计变更后更新滞后}——改一个组件，哪些 FMEA 行受牵连？
表格回答不了，只能人工全面复查，一次变更的 FMEA 更新要花\textbf{两周}；
\textbf{其三，知识不可复用}——新机型的液压泵故障模式与上一型号大同小异，
但上一型号的知识锁在表格与老分析师的记忆里，新项目几乎从零开始，
整个分析周期长达\textbf{六个月}。
读者应该已经嗅到熟悉的味道：这又是第 1 章的知识熵增，
而药方与前三个案例一脉相承——把表格升格为本体。

\subsection{本体骨架：概念与关系}

FMEA 本体的骨架在仓库文件 \texttt{aerospace-fmea.txt} 中列得干净利落，
七个核心类加六条关系，整理成一张表：

\begin{center}
\begin{tabularx}{\linewidth}{@{}lXX@{}}
\toprule
\textbf{构件} & \textbf{含义} & \textbf{示例或值域说明} \\
\midrule
System / Component & 系统与组件，沿 \texttt{partOf} 组成整机分解结构 & 液压系统；泵、阀门、管路 \\
FailureMode & 故障模式，「坏法」的类型化 & 泄漏、堵塞、断裂 \\
Effect & 影响，分局部、系统、飞行三层 & 流量减少 \(\rightarrow\) 压力丧失 \(\rightarrow\) 控制面失效 \\
Cause / Detection / MitigationAction & 原因、探测方法、缓解措施 & 油液污染；压力传感器；冗余系统 \\
hasFailureMode / hasCause & 组件 \(\rightarrow\) 故障模式（一对多）；故障模式 \(\rightarrow\) 原因 & 表格的「行」由此解构成图 \\
causes / propagatesTo & 故障模式 \(\rightarrow\) 影响；影响逐层传播（可传递） & 传播链的承重梁 \\
detectedBy / mitigatedBy & 故障模式 \(\rightarrow\) 探测；影响 \(\rightarrow\) 缓解 & 兜底措施挂接点 \\
\bottomrule
\end{tabularx}
\end{center}

骨架里藏着两个决定成败的设计决策。

\textbf{决策一：影响（Effect）是一等公民类，而不是故障模式的文本属性。}
表格时代「系统影响」只是一格文字；本体把每个影响建成\textbf{个体}，
于是「液压系统压力丧失」全库只有一个身份，
多个故障模式可以 \texttt{causes} 它，它再 \texttt{propagatesTo} 多个下游影响——
影响从孤立的描述变成传播网络中的节点，链式推理才有了落脚点。

\textbf{决策二：影响分三层并以 \texttt{propagatesTo} 串联。}
局部影响（组件层）、系统影响（分系统层）、飞行影响（整机层）
对应着适航分析的标准视角；
\texttt{propagatesTo} 声明为可传递，「局部影响最终牵动哪些飞行影响」
交给第 5 章的传递推理自动补全——这与制造案例 \texttt{precedes} 的手法如出一辙，
工程师只录相邻一跳，长链由推理器接力。

\subsection{影响传播链精讲}

骨架之上是规则。仓库文件给出三条 SWRL 规则，恰好覆盖三类推理任务。

\begin{codebox}
\codeline{\textcolor{cmtgray}{\#\ ==============================}}
\codeline{\textcolor{cmtgray}{\#\ SWRL推理规则（自动影响推导）}}
\codeline{\textcolor{cmtgray}{\#\ ==============================}}
\codeblank{}
\codeline{\textcolor{cmtgray}{\#\ 规则：液压泵故障\ \(\rightarrow\)\ 液压系统压力丧失}}
\codeline{Component(?c)\ \(\land\)\ hasFailureMode(?c,\ ?fm)\ \(\land\)}
\codeline{HydraulicPump(?c)\ \(\land\)\ FailureMode\_Blockage(?fm)}
\codeline{\hspace*{4\ttcw}\(\rightarrow\)\ SystemEffect\_PressureLoss(?fm)}
\codeblank{}
\codeline{\textcolor{cmtgray}{\#\ 规则：液压系统压力丧失\ \(\rightarrow\)\ 飞行控制面失效}}
\codeline{SystemEffect\_PressureLoss(?se)\ \(\land\)\ affects(?se,\ ?cs)\ \(\land\)}
\codeline{ControlSurface(?cs)}
\codeline{\hspace*{4\ttcw}\(\rightarrow\)\ FlightEffect\_ControlLoss(?se)}
\codeblank{}
\codeline{\textcolor{cmtgray}{\#\ 规则：设计变更影响分析（可追溯性）}}
\codeline{DesignChange(?dc)\ \(\land\)\ affects(?dc,\ ?component)\ \(\land\)}
\codeline{hasFailureMode(?component,\ ?fm)}
\codeline{\hspace*{4\ttcw}\(\rightarrow\)\ requiresFMEAUpdate(?dc,\ ?fm)}
\codeblank{}
\end{codebox}

\snipsource{ch06-applications/examples/aerospace-fmea.txt}

\textbf{第一条规则：从组件故障到系统影响。}
规则体绑定一个组件 \texttt{?c} 及其故障模式 \texttt{?fm}，
再用两个类原子收紧条件——组件得是液压泵（\texttt{HydraulicPump}），
故障模式得是堵塞（\texttt{FailureMode\_Blockage}）；
满足即把 \texttt{?fm} 归入「系统影响：压力丧失」类。
这是领域知识的最小颗粒：「液压泵堵塞会导致系统失压」，
原本写在老分析师的经验里，现在写成了机器可执行的一行。

\textbf{第二条规则：从系统影响到飞行影响。}
压力丧失 \texttt{?se} 若波及（\texttt{affects}）某个飞行控制面 \texttt{?cs}，
则升级为「飞行影响：控制能力丧失」。
妙处在于\textbf{两条规则首尾相扣}：
第一条贴上的「压力丧失」标签正是第二条的触发前提，
第 5 章讲过的前向链推理在此自动接力——
分析师没有写过「泵堵塞会威胁飞行控制」这句话，
它是推理器把两条局部知识\textbf{级联}出来的新结论。
故障传播链越长，这种级联的威力越大：
人脑推演链条会断，规则引擎不会。

\textbf{第三条规则：设计变更的可追溯性。}
设计变更 \texttt{?dc} 影响到组件 \texttt{?component}，
该组件挂着故障模式 \texttt{?fm}，
则断言 \texttt{requiresFMEAUpdate(?dc, ?fm)}——
「这次变更需要复查哪些 FMEA 条目」从两周的人工排查变成一次查询。
仓库文件记载这项能力把设计变更后的 FMEA 更新周期从\textbf{两周压缩到两天}，
而且给出的是\textbf{完备清单}：推理器沿着 \texttt{affects} 与
\texttt{hasFailureMode} 的边走，到得了的全到，不会像人一样凭印象挑着查。

三条规则放回 6.1 的模式语言里看：前两条是\textbf{决策支持}
（替分析师推演后果，结论自带推理链），
第三条是\textbf{合规审计}的基建（变更与分析活动全程留痕）。
一个规则库同时供养两种模式，这正是 FMEA 案例被称为「双模并重」的原因。

\subsection{液压系统实例走查}

把镜头拉近到一个具体故障，看本体里的一次完整推演。

\begin{codebox}
\codeline{\textcolor{cmtgray}{\#\ ==============================}}
\codeline{\textcolor{cmtgray}{\#\ 飞机液压系统FMEA示例}}
\codeline{\textcolor{cmtgray}{\#\ ==============================}}
\codeline{\textcolor{cmtgray}{\#\ 组件：液压泵(HydraulicPump\_01)}}
\codeline{\textcolor{cmtgray}{\#\ 故障模式：堵塞(Blockage)}}
\codeline{\textcolor{cmtgray}{\#\ 原因：液压油污染}}
\codeline{\textcolor{cmtgray}{\#\ 局部影响：泵输出流量减少}}
\codeline{\textcolor{cmtgray}{\#\ 系统影响：液压系统压力丧失}}
\codeline{\textcolor{cmtgray}{\#\ 飞行影响：飞行控制面响应迟滞，严重时失控}}
\codeline{\textcolor{cmtgray}{\#\ 探测方法：压力传感器、流量监测}}
\codeline{\textcolor{cmtgray}{\#\ 缓解措施：冗余液压系统、定期液压油更换}}
\end{codebox}

\snipsource{ch06-applications/examples/aerospace-fmea.txt}

这份示例是一条教科书式的 FMEA 记录，逐行翻译成本体断言：
个体 \texttt{HydraulicPump\_01} 属于液压泵类，
经 \texttt{hasFailureMode} 挂接堵塞故障模式个体；
故障模式经 \texttt{hasCause} 指向「液压油污染」，
经 \texttt{detectedBy} 指向压力传感器与流量监测两个探测手段。
影响链三级跳：局部影响「泵输出流量减少」、
系统影响「液压系统压力丧失」、飞行影响「控制面响应迟滞，严重时失控」，
由 \texttt{propagatesTo} 逐级相连；
「冗余液压系统」与「定期换油」作为缓解措施经 \texttt{mitigatedBy} 挂在影响上。

录入的只有这些\textbf{相邻事实}；上一小节的规则随即开始工作——
堵塞触发第一条规则归入压力丧失，压力丧失波及控制面触发第二条规则升级为飞行影响。
分析师在评审会上展示的不再是一行行表格，而是一张可以\textbf{点开任何一环}的传播图：
每个结论旁挂着触发它的规则与事实，质疑任何一环都能当场追到源头。
对适航审定而言，这张图就是证据链——
审查方问「为什么认定该故障威胁飞行安全」，答案不是「专家判断」，
而是两条可检视的规则加一串可核对的断言。

还有一个容易被忽略的收益藏在缓解措施一栏：
「冗余液压系统」本身也是本体里的系统个体，
推理器可以继续追问——\textbf{冗余系统自己的故障模式是什么？
主备同因失效（共因故障）的路径存在吗？}
表格时代这类追问要靠分析师的警觉，图谱时代它是一条查询。

\subsection{知识复用与适航视角}

FMEA 本体最大的红利在\textbf{跨型号复用}。
液压泵的故障模式不随机型改变——堵塞、泄漏、轴承磨损，
上一型号验证过的故障知识对新型号大体成立。
本体形态下，这些知识是结构化资产：
新机型启动 FMEA 时先\textbf{继承}通用故障知识库
（组件类的典型故障模式、影响模板、探测与缓解手段），
分析师的工作从「从零编表」变成「核对差异、补充新增」。
仓库文件记载的总账：分析周期从\textbf{六个月}缩短到\textbf{六周}（压缩 75\%），
故障模式识别量反而\textbf{提升 40\%}——
快了反而更全，因为推理器穷举传播路径不知疲倦，
而人在六个月的表格马拉松里注定会漏。
「更快」与「更全」在人工范式下是一对矛盾，在知识复用范式下是同一件事的两面。

从适航视角看，这套系统的价值还要再加一层：\textbf{分析过程本身可审计}。
哪个版本的故障知识库、哪些规则、哪一版设计数据推出了哪些结论，
全部有据可查；设计更改后哪些分析被标记重做、是否完成，
\texttt{requiresFMEAUpdate} 的断言记录就是台账。
第 5 章说过推理结果应当随推理器与本体版本一起归档，
在适航场景这不是良好实践，而是审定要求。

教训同样要记一条：\textbf{故障知识库需要治理}。
跨型号复用意味着一处错误会被复制到所有新项目，
知识库的修改权限、评审流程、版本基线必须像源代码一样管理——
这把我们引向第 7 章的知识图谱管理与第 8 章的质量评估，
也解释了为什么成熟的航空企业把本体团队设在系统工程部门而不是 IT 部门：
知识资产的治理责任必须落在懂知识的人身上。

\section{跨案例分析：共性、ROI 与风险}

\subsection{四个案例同框比对}

四个案例分属四个行业，节拍从百毫秒到数月，把它们放进同一张表里，
共性立刻浮出水面：

\begin{center}
\begin{tabularx}{\linewidth}{@{}lXX@{}}
\toprule
\textbf{案例（领域）} & \textbf{知识源与推理类型} & \textbf{价值度量} \\
\midrule
制造调度 & 工艺文件、设备台账、调度员经验；SWRL 冲突规则加传递属性推理 & 冲突检测自动化（响应小于 100 毫秒），工艺推荐命中率 87\% \\
自动驾驶 & 交通规则、高精地图、感知数据流；实例分类加时序风险规则 & 场景图每 100 毫秒更新，风险判定可解释、可回放 \\
BIM 合规 & IFC 模型、规范条款库；派生分类加 SWRL 检查加 SPARQL 报告 & 人工 2 天 \(\rightarrow\) 自动 30 分钟，条款覆盖 100\% \\
航空 FMEA & 故障知识库、设计文档、适航要求；链式 SWRL 影响传播推理 & 周期 6 个月 \(\rightarrow\) 6 周，识别量 +40\%，变更更新 2 周 \(\rightarrow\) 2 天 \\
\bottomrule
\end{tabularx}
\end{center}

逐列细读这张表，能提炼出四条跨行业的共性规律。

\textbf{规律一：知识源永远是「文档加系统加人脑」三件套。}
四个案例无一例外——结构化数据躺在业务系统里，
显式规则写在文档与规范里，最值钱的判断经验存在专家头脑里。
本体工程的本质工作就是把三种形态的知识\textbf{汇流到同一种表示}，
这正是第 3 章知识获取方法要解决的问题。

\textbf{规律二：推理的主力形态是「分类加规则」。}
没有一个案例动用了第 5 章价目表上最昂贵的表达力；
反而是「派生类自动归类（防火门、外墙）加 SWRL 贴标签（冲突、风险、违例）」
这对朴素组合包打了天下。工业本体的甜蜜点在轻表达力、重数据吞吐一侧，
表达力是按需采购的，不是越多越好。

\textbf{规律三：推理器只暴露问题，不替人决策。}
冲突清单交给调度员、风险断言交给规划器、违例报告交给审查师、
传播链交给分析师——四个案例不约而同地把「判定」与「处置」分开。
这既是技术审慎（规则覆盖不了全部情境），更是组织智慧：
系统帮人而不是换人，落地阻力小一个数量级。

\textbf{规律四：价值度量都是「时间加覆盖率」双指标。}
快（2 天到 30 分钟、6 个月到 6 周）只是表层，
全（100\% 条款、识别量 +40\%）才是质变——
机器检查的真正优势不是速度，而是\textbf{不知疲倦的完备性}。
向管理层论证项目价值时，这两类指标都要量化采集，缺一不可。

还值得把四个案例按\textbf{节拍}排成一条谱系：
自动驾驶 100 毫秒、制造调度分钟级、BIM 检查以设计迭代为周期、FMEA 以型号研制为周期。
同一副「TBox 稳定、ABox 流动、规则贴标签」的骨架，
在不同节拍下呈现出不同的工程形态——
节拍越快，表达力裁剪与增量推理越关键（6.3 的预算思维）；
节拍越慢，知识治理与版本追溯越关键（6.5 的治理教训）。
评估自己的场景时，先找到它在这条谱系上的位置，
就知道该把工程力气优先花在哪一侧。

\subsection{ROI 三判据}

「我的场景值不值得上本体？」综合四个案例，投入产出比可以用三条判据预判，
满足越多、回报越确定：

\begin{itemize}
\item \textbf{判据一：决策依据是知识而非统计。}
  排产约束、交通规则、规范条款、故障机理——四个案例的判断依据
  都能被专家\textbf{明示成条文}。如果你的问题主要靠海量样本里的统计规律
  （比如点击率预估），主角应该是机器学习，本体最多做特征与解释的配角；
\item \textbf{判据二：知识与数据的变化频率差足够大。}
  TBox 月变、ABox 秒变，分层架构才有红利；
  若概念体系本身天天变（早期探索型业务），建模会沦为追赶，先等领域稳定下来；
\item \textbf{判据三：错误代价高且需要解释。}
  停机、事故、审定失败——违例漏检的损失远超建模投入，
  且「为什么这样判」必须能回答。如果错了无所谓、也没人问为什么，
  上推理器是高射炮打蚊子。
\end{itemize}

三条判据其实分别回答三个朴素问题：知识在哪里（判据一）、
架构有没有红利（判据二）、值不值得较真（判据三）。
立项评审时把它们当一票否决项用——三条全不满足的项目，再时髦也不要做。

\subsection{失败风险清单}

成功案例讲完，更要紧的是失败的方式。四个案例的复盘加上业界的常见教训，
凝成一张风险清单，每条附解药：

\begin{itemize}
\item \textbf{范围蔓延}：什么都想建模，本体长成无人能维护的巨兽。
  解药：CQ 当闸门，没有问题牵引的概念一律不进（6.2 的疫苗）；
\item \textbf{数据管道烂尾}：本体精美，ABox 没人喂，推理器空转。
  解药：立项先验数据可得性，数据接入与本体建模同步排期（6.2 教训一）；
\item \textbf{规则库膨胀}：规则越写越多、互相纠缠，没人敢改。
  解药：宁少勿滥，先有真实案例再写规则，每条规则有负责人与测试用例；
\item \textbf{表达力滥用}：把价目表上最贵的构造当玩具用，性能塌方在上线日。
  解药：按第 5 章价目表做表达力预算，实时场景裁剪到规则匹配级（6.3 的预算思维）；
\item \textbf{形式化越界}：把模糊条款、模糊经验强行写成规则，结论似是而非。
  解药：诚实分级，可形式化的交机器，剩下的留给人并留痕（6.4 的条款分级）；
\item \textbf{知识库失治}：人人可改或无人可改，错误知识随复用扩散。
  解药：像管源代码一样管本体——权限、评审、版本、基线（6.5 的治理教训）；
\item \textbf{组织错位}：本体团队挂在纯 IT 编制下，离领域专家太远。
  解药：知识资产的所有权交给业务侧，IT 提供平台与工程支持。
\end{itemize}

把这张清单倒过来读，就是一份立项尽职调查问卷：
每个风险都问一句「我们有对应的解药吗」，答不上来的就是项目的真实短板。
值得强调的是，这七类风险没有一类是逻辑或算法层面的——
本体项目极少死于推理器，几乎总是死于范围、数据、治理与组织这些「工程之外的工程」。

\section{迁移指南：把案例经验带回你的领域}

\subsection{九步迁移检查单}

读完四个案例，最该带走的不是任何一个领域的本体，
而是把方法迁移到\textbf{你的}领域的能力。
下面这份检查单把四个案例的共同路径压缩成九步，分三个阶段：

\textbf{阶段一：立项判断（第 1--3 步）}
\begin{enumerate}
\item \textbf{识别主导模式}：对照 6.1 的四种模式与适用判据，
  确定你的场景以哪种模式为主、哪种为辅——这决定第一期交付物是什么；
\item \textbf{过 ROI 三判据}：知识可明示吗？变化频率分层吗？错误代价高吗？
  三问皆否则就此打住，把这本书省下的时间用在别处；
\item \textbf{盘点知识源与数据源}：列出文档（工艺、规范、手册）、
  系统（MES、ERP、台账）、专家（谁、还在不在、有多少时间），
  对每个数据源标注「接入难度」——数据管道烂尾是头号死因，丑话先讲在前面。
\end{enumerate}

\textbf{阶段二：最小闭环（第 4--7 步）}
\begin{enumerate}
\item \textbf{写 5--8 条能力问题}：用第 3 章的方法，
  让一线人员说出「最希望系统替你回答的问题」，按价值排序取前几条；
\item \textbf{圈定最小概念集}：只为这几条 CQ 建 TBox——
  四个案例的顶层类都不超过十个，你的第一版也不该超过；
\item \textbf{写 2--3 条规则跑通闭环}：选最痛的一类违例或冲突，
  规则照抄「规则体描述坏组合、规则头贴标签」的骨架，
  用真实历史数据回放验证——历史上出过的事故能被检出，闭环才算通；
\item \textbf{用 CQ 验收}：每条 CQ 翻译成 SPARQL 跑通，答错的回去改模型，
  这是第 8 章评估方法的最简版。
\end{enumerate}

\textbf{阶段三：扩展运营（第 8--9 步）}
\begin{enumerate}
\item \textbf{接通数据管道}：把 ABox 从手工录入换成系统对接，
  推理从手动触发换成事件或周期触发，此时才谈得上「上线」；
\item \textbf{建立治理机制}：本体与规则的修改流程、版本基线、
  负责人与值班表——参照 6.6 风险清单逐项配解药，然后进入第 7 章的运营正题。
\end{enumerate}

\subsection{迷你示例：注塑车间走前三步}

检查单是抽象的，拿一个\textbf{虚构的注塑车间}走一遍前三步，看看每步产出什么。
设想这家工厂有几十台注塑机（按锁模力分大小机型）、几百副模具、
常用塑料粒子若干种（如 ABS、PC），订单由 ERP 下发，
排产靠一位老调度和一块白板，常见事故是
「模具装到锁模力不够的机台上」与「物料与模具不匹配导致批量缺陷」。

\textbf{第一步产出：主导模式判定。}
事故清单里的两类问题都是「坏组合」——模具对机台、物料对模具，
典型的\textbf{约束检查}场景；老调度的换模顺序经验则指向辅助的\textbf{决策支持}。
与 6.2 的机加工车间几乎同构，第一期交付物因此明确：排产冲突自动检测。

\textbf{第二步产出：ROI 三判据过堂。}
判据一：「多大锁模力配多大模具」「哪种粒子能走哪副模具」
都写得进表格、说得成条文——知识可明示，过；
判据二：机台与模具的能力参数年度才变，订单与排产分钟级在变——频率分层明显，过；
判据三：装错模具轻则停机换装、重则压坏模具（一副精密模具价值不菲），
且出了事要追「谁排的、为什么」——代价高且需解释，过。三判据全过，立项成立。

\textbf{第三步产出：知识源与数据源清单。}
文档类：模具台账（含适用机型、腔数、保养记录）、物料规格书；
系统类：ERP 订单接口（现成）、注塑机状态信号（部分老机台无数采，标注「接入难」）；
专家类：老调度一位（换模顺序与「哪台机干哪类活最稳」的手感在他脑中，
访谈窗口每周约两小时）。
清单里「老机台无数采」一条立刻预警了 6.6 风险清单的第二项——
第一期规则应避开依赖实时状态的判定，先做静态的能力匹配检查，
状态类规则等数采改造完成后再上。

三步走完，项目还没写一行公理，但已经知道\textbf{做什么、值不值、缺什么}。
接下来的第四步起，剧本就是 6.2 的重演：
为「哪些机台能装这副模具」「哪些排产存在锁模力不匹配」写 CQ，
建机台、模具、物料、任务四个核心类，
写下第一条规则——规则体绑定任务、模具与机台，
比较所需锁模力与机台额定锁模力，规则头贴上冲突标签。
你会发现 6.2 那三条规则的骨架原封不动地适用，换的只是名词。
\textbf{这就是模式的力量：案例不可复制，骨架可以。}

\begin{noteblock}
\textbf{第一期不要超过一个季度。}
四个案例的共同节奏是「小闭环先跑通、价值先可见」：
制造案例先上三条冲突规则，BIM 案例先做量化阈值类条款。
一个季度内交付不了可演示闭环的方案，多半是范围切大了——
回到检查单第 5 步，把概念集再砍一半。
\end{noteblock}

\section{本章小结}

本章把前五章备好的理论与工具送进了四个真实领域，沿途反复验证同一副骨架。
\textbf{模式层}：本体落地的四种应用模式——约束检查、语义集成、决策支持、合规审计，
各有通用结构与适用判据，真实项目通常一主一辅地叠加使用。
\textbf{案例层}：制造调度用三条 SWRL 规则让能力、状态、时间三类排产冲突无处可藏，
冲突响应小于 100 毫秒、工艺推荐命中率半年爬到 87\%；
自动驾驶把感知结果填成每 100 毫秒更新的场景图，
TTC 与行人危险区规则给出可回放、可审查的风险断言，
推理预算与增量更新是实时性的生命线；
BIM 合规检查以 IfcOWL 打通模型词汇与规范词汇，
量化条款沿五步流水线形式化，人工 2 天的审查变成自动 30 分钟、条款覆盖 100\%，
同时诚实划界——模糊条款仍归人工；
航空 FMEA 把故障表格升格为传播网络，两条规则级联出「泵堵塞威胁飞行控制」的新结论，
变更影响排查从 2 周到 2 天，分析周期从 6 个月到 6 周、识别量反升 40\%。
\textbf{方法层}：跨案例的四条共性规律、ROI 三判据、失败风险清单与九步迁移检查单，
是把案例经验搬回你自己领域的工具包。

四个案例的每一处关键机制都能在前几章找到出处：
CQ 与范围控制来自第 3 章，类与属性的表达来自第 4 章，
规则、传递属性与前向链来自第 5 章，分层与开放世界的底层逻辑来自第 2 章；
而它们暴露的新问题——数据管道、版本治理、质量评估——
正是第 7 章知识图谱管理与第 8 章评估方法的开场白。
第 9 章将以制造调度为蓝本，从零写出一个可单步调试的完整系统，
本章读出的骨架届时会变成你亲手敲下的代码。

给动手派读者一条练习路线：
先把四个仓库文件整篇通读一遍——正文为叙事顺畅而省略的属性清单与注释细节都在里面；
然后任选一个最贴近你行业的案例，把它的类层次与规则录入 Protégé，
造三五个个体、跑一遍推理，亲眼看看冲突或违例标签是怎么被贴上的；
最后用 6.7 的检查单对自己的领域做一次纸面推演，写出主导模式判定、
三判据结论与 5 条 CQ——这份一页纸的产出，就是你下一个本体项目的种子。

\subsection*{思考题}

\begin{enumerate}
\item 给下列三个场景判定主导应用模式并说明理由：
  电网检修计划的安全规程审查；医院处方系统的药品相互作用提示；
  铁路信号系统多供应商数据的统一接入。
  提示：对照 6.1 各模式的适用判据，注意「报告给谁看」往往暴露主导模式。
\item 为 6.2 的调度本体补一条 SWRL 规则：处于「维护中」状态的设备不得被指派任务。
  写出全部原子，并说明它与状态冲突规则的异同。
  提示：仿照状态冲突规则的骨架，只需更换状态值原子。
\item 能力冲突规则里的否定原子在开放世界假设下为什么危险？
  给出正文提到的两种工程绕法，并比较各自的维护成本。
  提示：回看 6.2.5——「没说能加工」与「不能加工」在 OWL 里不是一回事。
\item 自动驾驶的 TTC 规则为什么必须先用 \texttt{greaterThan} 判定自车更快，
  再做减法与除法？从业务语义与技术保护两个角度各说一条理由。
  提示：想一想速度差为零时除法会发生什么，以及内建原子的绑定顺序纪律。
\item 某节能条款的限值随气候分区与建筑类型二维变化。
  仿照 6.4.4 的参数化手法，设计所需的类、属性与规则骨架。
  提示：把「查表」建成 ABox 断言，规则里用变量把限值查出来再比较。
\item 若把 FMEA 本体的 \texttt{propagatesTo} 声明为传递属性，
  哪类能力问题直接受益？长传播链可能带来什么风险？
  提示：对照第 5 章传递推理的代价，想想几十级传播链上的结论还可不可信。
\item 用 6.6.2 的 ROI 三判据评估你熟悉的一个业务场景，
  给出三条判据的逐条结论与总体建议。
  提示：判据二最常被忽略——先画出该场景里「慢变知识」与「快变数据」各是什么。
\item 接着 6.7.2 的注塑车间走第四到第六步：写出 5 条 CQ、
  画出核心类草图，并写出「锁模力不匹配」检测规则的草稿。
  提示：6.2 三条规则的骨架原封不动适用，把材料换成模具与锁模力即可。
\end{enumerate}

\begin{noteblock}
\textbf{做题建议。}
每题末尾的提示只指方向，不给答案；
动笔前先回到对应小节把片段重读一遍，
第 2、5、8 题写完后强烈建议录入 Protégé 实际跑一跑——
规则「看着对」与「跑得对」之间的距离，正是本章想让你亲手量一次的。
\end{noteblock}

\subsection*{本章文件指引}

本章全部代码片段节选自下列四个仓库文件；
它们是案例的「原始档案」，注释里保有正文未及展开的细节，
建议对照正文整篇通读。

\begin{center}
\begin{tabularx}{\linewidth}{@{}>{\ttfamily\footnotesize}p{0.40\linewidth}X@{}}
\toprule
\normalfont\textbf{文件（ch06-applications/examples/）} & \textbf{内容与阅读建议} \\
\midrule
manufacturing-scheduling.txt & 调度本体知识流架构、五条 CQ、完整类层次与属性清单、三条 SWRL 冲突规则与量化效果；配合 6.2 与第 9 章实战阅读 \\
autonomous-driving.txt & 场景图 100 毫秒更新流程、参与者与道路类层次、空间关系属性、TTC 与行人危险区规则；配合 6.3 阅读 \\
bim-ontology.txt & 两处断层、IfcOWL 方案四步、核心类 Manchester 语法、防火与节能检查规则、SPARQL 合规查询；配合 6.4 与第 4 章查询语法阅读 \\
aerospace-fmea.txt & FMEA 七类六关系、三条影响传播与变更追溯规则、液压系统完整示例、量化效果；配合 6.5 与第 5 章规则机制阅读 \\
\bottomrule
\end{tabularx}
\end{center}
